\begin{abstract}
This manuscript is the first real attempt to turn the repository into one paper rather than a stack of local reports. It consolidates the full ARC analysis line in the workspace: LLM psychometrics, human testing, non-LLM comparisons, validated Python-solver complexity, cross-ARC latent linkage, efficiency analysis, and the external small-model ARC literature. The goal is not to restate each report independently, but to reconstruct the shared methodology, the overlap structure, the strongest findings, and the places where the current evidence is still fragile. The resulting picture is coherent but not simplistic. Humans and LLMs do appear to share a real ARC difficulty axis, but that shared structure is incomplete. The strongest and most stable result in the repository is that validated solver structure is a strong predictor of LLM task difficulty and a much weaker predictor of human difficulty. Human burden looks more tied to duration, search, and pair-level heterogeneity than to the size of the final successful solver. Current non-LLM systems are the weakest part of the empirical picture if the target is broad human-like replication: they provide complementary successes and some residual signal, but they do not yet reproduce the human difficulty profile as well as the LLM consensus does. At the same time, the external architecture papers suggest that non-LLM progress converges on a loop built from strong inductive bias, iterative refinement, and test-time candidate selection. The strongest unified thesis is therefore not that one paradigm has solved ARC ``the human way.'' It is that ARC contains a shared abstraction factor plus source-specific burdens, with LLMs appearing especially structure-sensitive, humans appearing more search- and interaction-sensitive, and current non-LLM systems occupying an interesting but still underpowered middle position.
\end{abstract}

\tableofcontents
\clearpage
